\section{Introduction}

Late interaction models, also referred to as multi-vector or ColBERT~\cite{khattab2020colbert} models, have gained popularity thanks to their out-of-domain~\cite{santhanam2022colbertv2}, long-context~\cite{chaffin2025gte}, and reason-intensive~\cite{chaffin2025reason} retrieval capabilities. Despite these clear advantages, they are less studied than their single-vector counterparts. Thus, the state-of-the-art ColBERT models such as GTE-ModernColBERT~\cite{chaffin2025gte} and ColBERT-small~\cite{clavie2025colbert} are built by performing a small Knowledge Distillation (KD) phase on top of a strong pre-trained single-vector (dense) model. 
Even mxbai-edge-colbert-v0~\cite{takehi2025fantastic}, a recent study that details a recipe to train very strong ColBERT models from scratch (that is, MLM-pretained Ettin~\cite{weller2025seq} models) performed all the early stages in a single-vector setting and only used multi-vector in the latest stage, knowledge distillation.

Training a retrieval model typically involves up to three phases (illustrated in Figure~\ref{2602:fig:training_phase}): (1) \emph{unsupervised contrastive pre-training}, a large-scale contrastive phase relying on in-batch negatives; (2) \emph{supervised contrastive fine-tuning}, which refines the model using mined hard negatives; and (3) \emph{Knowledge Distillation} (KD), where a strong teacher's relevance scores guide the student via KL divergence. 
Although applying only KD yields strong models and allows us to leverage various strong dense models, it treats the multi-vector setting as an afterthought. This approach could be suboptimal compared to performing contrastive steps in the multi-vector setting before the knowledge distillation step. 

In this paper, we investigate which training steps are necessary for optimal ColBERT performance. Specifically, we ask two questions: is knowledge distillation alone sufficient to transfer dense model quality to the multi-vector setting? And if not, can a supervised contrastive phase before KD close the gap without the costly unsupervised pre-training? To answer these, we compare models trained with (a) the knowledge distillation step alone, (b) the supervised contrastive and knowledge distillation steps, and (c) all steps, including the large-scale unsupervised contrastive step; see Figure~\ref{2602:fig:pipelines} for details. Our findings show that KD alone is insufficient for optimal performance, but that adding a supervised phase beforehand yields competitive results without the most expensive unsupervised phase, particularly when aligning the fine-tuning setup with the pre-training one. Indeed, we find that aligning the use of prompts is crucial when repurposing existing models, with misalignment leading to significant performance degradation The resulting model, \emph{ColBERT-Zero}, trained only on public data, outperforms both the state-of-the-art GTE-ModernColBERT as well as its base model, GTE-ModernBERT, trained on closed and much stronger data, and sets new state-of-the-art for its size. To support reproducibility and further exploration, we release \href{https://huggingface.co/collections/lightonai/colbert-zero}{all models}, including intermediate checkpoints with various training parameters, as well as the \href{https://github.com/lightonai/pylate/tree/main/examples/train/ColBERT-zero/}{full training scripts}.
